% docs/ecfi_minipaper_trace.tex
\documentclass[11pt]{article}

\usepackage[a4paper,margin=2.6cm]{geometry}
\usepackage{amsmath,amssymb}
\usepackage{booktabs}
\usepackage{graphicx}
\usepackage{hyperref}
\usepackage{enumitem}

\hypersetup{
  colorlinks=true,
  linkcolor=blue,
  citecolor=blue,
  urlcolor=blue
}

% ---- Notation helpers (avoid N vs \mathcal{N} collision) ----
\newcommand{\R}{\mathbb{R}}
\newcommand{\Neigh}{\mathcal{N}} % neighbor set
\newcommand{\nov}{\nu}          % novelty
\newcommand{\varph}{\varphi}    % prefer \varphi
\newcommand{\norm}[1]{\left\lVert #1 \right\rVert}

\title{Tracing a Minimal E-CFI Demo: From Architecture to Prototype, Diagnostics, and Step H$'$ Tuning}
\author{Internal mini-paper (repo trace)}
\date{\today}

\begin{document}
\maketitle

\begin{abstract}
This note documents, in a traceable way, the first runnable prototype of an Embodied Collective Free Improvisation (E-CFI)
swarm demo and the first diagnostic cycle obtained from simulation logs.
The prototype instantiates: (i) an embodied self-assembly substrate (join/detach and component tracking), (ii) event-driven oscillators
with merge-triggered phase synchronization, and (iii) a two-time-scale cognitive layer inspired by the collective free improvisation (CFI)
template (fast intention, slow objective), coupled to intrinsic signals (novelty, boreness, and load proxies).
We report a characteristic failure mode of the initial parameterization---rapid collapse into a single assembly component and permanent
phase lock-in---and we propose a minimal Step H$'$ tuning (plus an optional soft-sync rule) aimed at producing observable plateaus and
reconfiguration phases, as required by the vision paper.
\end{abstract}

\section{Prototype scope and implementation snapshot}
\subsection{What is implemented (Step G/H)}
The current prototype implements the following modules:

\begin{enumerate}[leftmargin=2.2em]
\item \textbf{Embodiment (assembly graph).} Agents move in a 2D torus. An \emph{assembly graph} $G_A(t)$ is formed by slot-limited contact joins
within radius $R_C$ and broken by probabilistic detachments after a minimum dwell time $T_{\min}$.
Connected components of $G_A(t)$ define \emph{assemblies} (clusters), whose sizes and counts are logged.
\item \textbf{Oscillators.} Each agent carries a phase $\varph_k(t)\in[0,1)$ with free-running drift $\dot{\varph}=\omega_0$ and optional noise.
When a join merges two previously disconnected components, phases are synchronized within the merged component (initially via \emph{hard sync}).
\item \textbf{CFI-inspired cognition.} Each agent computes a 3D local feature vector
$x_k(t)=(\deg_k(t),\ \text{cluster\_size}_k(t),\ \text{local\_coherence}_k(t))$.
A fast intention state $\omega_k(t)\in\R^3$ is updated by a saturated dynamical rule
\[
\frac{d\omega_k}{dt}=\alpha x_k(t) + \beta_k(t)\sum_{\ell\in\Neigh_k(t)}x_\ell(t) - g\norm{\omega_k(t)}^2\omega_k(t),
\]
while the slow objective is encoded by mixture weights $\Omega_k(t)=(w^I_k,w^C_k,w^N_k)$ controlling imitation/contrast/independence.
\item \textbf{Intrinsic constraints.} Novelty $\nov_k(t)$ is defined as the distance between successive EMAs of $x_k(t)$;
boreness $b_k(t)$ is a leaky integrator of novelty deficits, and load $c_k(t)$ is a proxy combining neighborhood size, assembly degree,
event rate, and novelty.
When $b_k>b_{\max}$ or $c_k>c_{\max}$, objective weights shift toward independence.
\item \textbf{Policy.} Motion uses a steering law combining alignment, cohesion, and separation, modulated by $\Omega_k$ and $\norm{\omega_k}$.
\end{enumerate}

\subsection{Reproducibility contract (how to rerun)}
A run is determined by a YAML config and a seed. The runner writes:
\texttt{outputs/run/config\_resolved.json}, \texttt{global.csv}, \texttt{agents.csv}, and diagnostic plots.
In the repository, the basic workflow is:

\begin{itemize}[leftmargin=2.2em]
\item run: \texttt{python -m ecfi.run --config configs/base.yaml --seed 0}
\item plot: \texttt{python -m ecfi.viz.plot\_run --run\_dir outputs/run}
\end{itemize}

\section{Observed behavior from the first diagnostic cycle}
Figures~\ref{fig:clusters}--\ref{fig:intrinsic} report the three core diagnostic plots obtained from the initial parameterization
(\texttt{configs/base.yaml}) with $N=50$ agents and $3000$ steps.

\begin{figure}[t]
  \centering
  \includegraphics[width=0.92\linewidth]{figures/plot_clusters.png}
  \caption{\textbf{Assembly regime indicators.}
  The number of connected components $n_{\text{clusters}}(t)$ drops quickly and stabilizes at $1$,
  while the mean cluster size grows to $50$ and remains there. This indicates a rapid transition to a single
  connected assembly with no subsequent fragmentation.}
  \label{fig:clusters}
\end{figure}

\begin{figure}[t]
  \centering
  \includegraphics[width=0.92\linewidth]{figures/plot_coherence.png}
  \caption{\textbf{Global phase coherence.}
  The order parameter saturates to $1$ and remains locked, i.e., complete global synchronization.
  Given the merge-triggered phase synchronization rule, this is consistent with the emergence of a single
  connected assembly coupled to a strong (hard) synchronization mechanism.}
  \label{fig:coherence}
\end{figure}

\begin{figure}[t]
  \centering
  \includegraphics[width=0.92\linewidth]{figures/plot_intrinsic.png}
  \caption{\textbf{Intrinsic signals (mean novelty, boreness, load).}
  Load remains in a relatively high band (driven by neighborhood and assembly degrees), while mean novelty exhibits
  short transients early on. Boreness appears nearly flat at the plotted scale, which matches the intended update rule:
  when novelty is near zero, boreness tends to a small value on the order of $\eta$ (hence it may not exceed $b_{\max}$
  for the initial thresholds).}
  \label{fig:intrinsic}
\end{figure}

\section{Diagnosis: why the initial parameterization locks in}
The diagnostic plots support the following causal reading.

\begin{enumerate}[leftmargin=2.2em]
\item \textbf{Assembly becomes globally connected.}
The join rule and steering policy produce a connectivity cascade, after which $G_A(t)$ has a single connected component.
\item \textbf{Hard phase synchronization creates an absorbing state.}
When a join merges components, the prototype uses hard sync:
$\varph_j \leftarrow \varph_{\text{anchor}}$ for all $j$ in the merged component.
Once connectivity reaches one component, repeated merges (or early merges) quickly drive the entire population to a shared phase,
after which global coherence remains at $1$.
\item \textbf{Boreness thresholds are not triggered.}
With the current definition
$b_k(t)=(1-\lambda)b_k(t-\Delta)+\lambda\max(0,\eta-\nov_k(t))$,
if $\nov_k(t)\approx 0$, then $b_k(t)$ tends to $\eta$ (not to an unbounded growth).
Thus, if $b_{\max}\gg \eta$, the objective update toward independence never activates.
\end{enumerate}

In short: the current configuration generates a strong, quickly absorbing regime (single assembly + full phase lock-in),
which prevents the demo from exhibiting alternating plateaus and reconfiguration phases.

\section{Step H$'$ tuning: minimal modifications to induce regime switching}
We propose a minimal Step H$'$ tuning that targets the three points above, without changing the architecture.

\subsection{Parameter-level tuning (config only)}
The following changes are sufficient to (i) increase fragmentation propensity, (ii) make boreness-triggered objective shifts operative,
and (iii) reduce unconditional cohesion.

\begin{center}
\begin{tabular}{@{}lll@{}}
\toprule
\textbf{Block} & \textbf{Baseline (H)} & \textbf{Tuned (H$'$)} \\
\midrule
assembly & $T_{\min}=40$, $p_{\text{detach}}=0.001$ & $T_{\min}=20$, $p_{\text{detach}}=0.008$ \\
constraints & $\eta=0.02$, $b_{\max}=1.0$, $c_{\max}=6.0$ & $\eta=0.05$, $b_{\max}=0.03$, $c_{\max}=4.3$ \\
objective & init $(0.70,0.20,0.10)$ & init $(0.60,0.15,0.25)$; faster update \\
policy & stronger cohesion & reduced cohesion; stronger separation \\
\bottomrule
\end{tabular}
\end{center}

Operationally, H$'$ should yield:
(i) clusters that form but are more likely to detach,
(ii) intrinsic triggers that shift weights toward independence,
(iii) oscillatory behavior in $n_{\text{clusters}}(t)$ and nontrivial coherence trajectories.

\subsection{Optional rule-level tuning: soft sync instead of hard sync}
If global coherence remains stuck at $1$ even after parameter tuning, we replace hard sync with a \emph{soft} synchronization step,
applied only on merge events:
\[
\varph_n \leftarrow (1-\rho)\varph_n + \rho\,\varph_{\text{anchor}} \pmod{1},
\qquad \rho\in(0,1).
\]
This preserves the embodied-merge semantics (a merge tends to synchronize) while avoiding an immediate absorbing lock-in.

\section{Traceability within the repository}
To keep a clean trace of this development cycle, we recommend:

\begin{enumerate}[leftmargin=2.2em]
\item Keep baseline configs and tuned configs separately:
\texttt{configs/base.yaml} (H) vs.\ \texttt{configs/regime\_switch.yaml} (H$'$).
\item Store figures in \texttt{docs/figures/} and include them in this note.
\item Record the git commit hash of each milestone in a short \texttt{docs/TRACE.md} file (G, H, H$'$),
and include the command outputs or the resolved configs saved under \texttt{outputs/}.
\end{enumerate}

\section{References}
\begin{thebibliography}{10}

\bibitem{canonne_garnier_mcm2011}
Cl{\'e}ment Canonne and Nicolas Garnier.
\newblock A Model for Collective Free Improvisation.
\newblock In \emph{Mathematics and Computation in Music (MCM), Third International Conference},
  \emph{Lecture Notes in Computer Science}, vol.~6726, pp.~29--41. Springer, 2011.
\newblock DOI: \href{https://doi.org/10.1007/978-3-642-21590-2_3}{10.1007/978-3-642-21590-2\_3}.

\bibitem{lucas_acsos2024_embodied}
Pedro Lucas, Alexander Szorkovszky, Stefano Fasciani, and Kyrre Glette.
\newblock Self-Assembly and Synchronization: Crafting Music with Multi-Agent Embodied Oscillators.
\newblock In \emph{Proceedings of the IEEE International Conference on Autonomic Computing and Self-Organizing Systems (ACSOS)},
  pp.~145--150. IEEE, 2024.
\newblock DOI: \href{https://doi.org/10.1109/ACSOS61780.2024.00034}{10.1109/ACSOS61780.2024.00034}.

\end{thebibliography}

\end{document}